\section{Reproducible Python implementation}
\label{sec:expanded-software}
The companion open-source \texttt{trendestimation} Python library
(version \texttt{0.1.0a0} at the evaluated snapshot) is itself
a practical outcome: it exposes Guerrero-inspired \emph{smoothness}
as a controllable coordinate while retaining the conventional
penalty $\lambda$ for compatibility. We distinguish this
software design and numerical implementation from the already
known statistical PLS method and classical selectors.
We do not claim that this is the first software package for
Whittaker--Henderson smoothing.

\subsection{Public API and reproducible usage}
A reusable \texttt{TrendSmoother(n, order=d)} object constructs
the difference operator and spectral decomposition once.
The user may prescribe a normalized target $s$, invert it to
a penalty and then fit; alternatively, a criterion selects
the penalty on either coordinate. The following minimal
example shows both workflows and the values returned:
\begin{quote}
\begin{verbatim}
from trendestimation import TrendSmoother, select_lambda

smoother = TrendSmoother(len(y), order=2)

# Prescribe normalized smoothness
lam = smoother.lambda_for_smoothness(0.86)
controlled = smoother.fit(y, lambda_=lam, mu=0.0)
print(controlled.smoothness)

# Or select an existing criterion in S coordinates
chosen = select_lambda(
    smoother, y, criterion="gcv", mu=0.0,
    search_space="smoothness",
)
trend = smoother.fit(
    y, lambda_=chosen.lambda_, mu=0.0,
).trend
print(chosen.lambda_, chosen.smoothness)
\end{verbatim}
\end{quote}
The criterion itself is unchanged; \texttt{search\_space}
chooses the numerical coordinate of its candidate search.
The returned \texttt{lambda\_} can be passed to
\texttt{fit}, while \texttt{smoothness} reports the
trace-normalized attained value. The implementation
also returns the evaluated penalty grid, scores, and
requested coordinate grid for auditing.
The current default public \texttt{TrendSmoother} uses
the reusable spectral backend; the benchmarked dense
and banded variants are not all automatic public
defaults selected by a benchmark ranking.

\subsection{Backend differences and caching}
The seven numerical implementations evaluated in the
large factorial are fresh dense solves, cached dense
solves, general banded solves, dense Cholesky,
banded Cholesky, symmetric eigendecomposition of
$Q_d$, and full SVD of $D_d$.
These are distinct linear algebra backends, \emph{not}
seven statistically different smoothers.

The full SVD expresses the penalty eigenvalues through
the squared singular values of $D_d$.
Once $y$ and a fixed $\mu$ have been projected,
evaluating traces and degrees of freedom for an additional
$\lambda$ costs only a sum over eigenvalues.
For a full reconstructed curve, the transformed spectral
coefficients must also be mapped back to $n$ coordinates.
Therefore a single EDF evaluation, one full fitted trend,
an entire CV grid including leverages, and multiple
independent searches are different timing problems.
A banded $I+\lambda Q_d$ admits sparse-band storage
and specialized factorization, often attractive for
single fits or larger $n$ \cite{weinert2007wh}.

\begin{table}[tbp]
\centering\footnotesize
\caption{Computational activities and what must be charged.
The table is methodological rather than a measured timing
ranking. $M$ is the penalty-candidate count and $R$
the number of series. Reused work must be charged
once in an end-to-end comparison.}
\label{tab:expanded-costs}
\begin{tabular}{p{.19\linewidth}p{.34\linewidth}p{.34\linewidth}}
\toprule
Task & Preprocessing and repeated work & Fair comparison criterion\\
\midrule
One fixed trend
& Build $Q_d$ or its band; solve for one $\lambda$
& Fit and setup time, agreement in $\widehat\tau$\\
One scalar criterion
& Compute fitted residuals and required EDF/leverage or likelihood diagnostics
& Score validity and total evaluation time\\
Full $M$-point search
& Prepare backend; evaluate $M$ objective scores
& Identical candidates, all scores valid\\
$R$ reused searches
& Share spectrum, candidate inversions or factorization
& Charge common preparation once\\
$S^{-1}(s)$ only
& Reuse prepared positive eigenvalues; iterate scalar root
& Tail error, failure rate and microseconds\\
\bottomrule
\end{tabular}
\end{table}

\subsection{Grid reuse and algorithmic workflow}
Figure~\ref{fig:expanded-schematic} outlines the
computational workflow. Given $(n,d)$, spectral
preprocessing depends on the \emph{operator}, not
on the particular series. A common grid
$s_1,\ldots,s_M$ can be inverted once,
cached as $\lambda_1,\ldots,\lambda_M$,
and shared among selectors and observed series.
That is different from computing a fresh inverse
for each criterion evaluation. If the entire grid
is used only once, its inversion costs must be
included in the end-to-end timing. For a
separately spaced log-$\lambda$ grid, the
candidate points differ, so the selected fit may
also differ.

\begin{figure}[tbp]
\centering
\setlength{\fboxsep}{8pt}
\fbox{\begin{minipage}{.91\linewidth}
\centering\small
{\bfseries Data and difference order $(y,n,d,\mu)$}\\[3pt]
$\Downarrow$\\[3pt]
{\bfseries Prepare $D_d$, spectral values or banded factorization}\\[3pt]
$\Downarrow$\\[3pt]
\begin{tabular}{p{.39\linewidth}|p{.39\linewidth}}
\textbf{Controlled $s$} & \textbf{Criterion selection}\\
$s\mapsto\Lambda_d(s)$ & $s$ grid or $\lambda$ grid\\
One safeguarded inverse & Optional shared inverse grid\\
\end{tabular}\\[5pt]
$\Downarrow$\\[3pt]
{\bfseries Same PLS fit at selected $\lambda$:
$(\widehat\tau,\lambda,S,\nu_\lambda)$}\\[3pt]
$\Downarrow$\\[3pt]
{\bfseries Save raw scores, selections, machine metadata,
numeric validity, and plots independently}
\end{minipage}}
\caption{Reusable software path for controlled smoothness
and criterion-based selection. Native-$S$ and once-mapped
$S\to\lambda$ grids are equivalent when their candidates
match; a separate log-penalty grid generally is not.
Each panel of the empirical application reads previously
saved selections without repeating benchmark timings.}
\label{fig:expanded-schematic}
\end{figure}

\subsection{Design for independent reproduction}
The repository separates numerical data production,
statistical analysis and figure rendering. The
80,640-trial Windows CSV and environment file, the
seed-paired solver records, and the scalar root-finding
results have been committed independently of all
PDF/SVG/PNG outputs. The plotting scripts verify
SHA-256 digests of the numeric archive and reconstruct
curves at \emph{stored} $(\lambda,\mu)$.
Thus a journal reader may change the color palette or
replace a figure layout without rerunning data-generating
simulations or hardware-sensitive timing experiments.
The independent original draft and complete experiment
matrices remain available as supplementary material.
